\documentclass[12pt]{article}
\usepackage{amsmath, amssymb, geometry}
\usepackage{array}
\usepackage{booktabs}

\geometry{margin=1in}

\title{Modeling Project: Decaying Average Grading System}
\author{}
\date{}

\begin{document}

\maketitle

\section*{Background}

Some grading systems use a \textbf{decaying average} so that more recent assignments count more than earlier ones. One common model is:

\[
G_n = 0.6 S_n + 0.4 G_{n-1}
\]

where:
\begin{itemize}
    \item $S_n$ = score on the newest assignment
    \item $G_n$ = updated grade
    \item $G_{n-1}$ = previous grade
\end{itemize}

\section*{Part 1: Understanding the Model}

\begin{enumerate}
    \item Suppose a student's first score is 70. What is their first decaying average?
    \vspace{1.5cm}
    
    \item If the next score is 80, compute the new grade.
    \vspace{1.5cm}
    
    \item If the next score is 90, compute the updated grade.
    \vspace{1.5cm}
    
    \item In your own words:
    \begin{itemize}
        \item Why might a school use a decaying average system?
        \item What does this system assume about how students learn?
    \end{itemize}
    \vspace{2cm}
\end{enumerate}

\section*{Part 2: Testing Learning Patterns}

For each dataset:
\begin{itemize}
    \item Compute the decaying average after each assignment
    \item Graph the scores and the decaying average
    \item Describe what the grading system concludes
\end{itemize}

\subsection*{Data Set A: Consistent Improvement}

\begin{center}
\begin{tabular}{cc}
\toprule
Assignment & Score \\
\midrule
1 & 55 \\
2 & 65 \\
3 & 75 \\
4 & 85 \\
5 & 92 \\
\bottomrule
\end{tabular}
\end{center}

\vspace{1cm}

\subsection*{Data Set B: Early Success, Later Struggle}

\begin{center}
\begin{tabular}{cc}
\toprule
Assignment & Score \\
\midrule
1 & 95 \\
2 & 90 \\
3 & 82 \\
4 & 70 \\
5 & 60 \\
\bottomrule
\end{tabular}
\end{center}

\vspace{1cm}

\subsection*{Data Set C: Fluctuating Performance}

\begin{center}
\begin{tabular}{cc}
\toprule
Assignment & Score \\
\midrule
1 & 90 \\
2 & 65 \\
3 & 88 \\
4 & 70 \\
5 & 92 \\
\bottomrule
\end{tabular}
\end{center}

\vspace{1cm}

\section*{Part 3: A Hidden Problem}

A class studies three units:
\begin{itemize}
    \item Linear Functions
    \item Quadratic Functions
    \item Exponential Functions
\end{itemize}

\subsection*{Student D}

\begin{center}
\begin{tabular}{ccc}
\toprule
Test & Topic & Score \\
\midrule
1 & Linear & 95 \\
2 & Quadratic & 96 \\
3 & Exponential & 55 \\
\bottomrule
\end{tabular}
\end{center}

\begin{enumerate}
    \item What does the decaying average say about this student?
    \vspace{1.5cm}
    
    \item What does the raw data suggest?
    \vspace{1.5cm}
    
    \item What information might the grading system miss?
    \vspace{1.5cm}
\end{enumerate}

\subsection*{Student E}

\begin{center}
\begin{tabular}{ccc}
\toprule
Test & Topic & Score \\
\midrule
1 & Linear & 55 \\
2 & Quadratic & 95 \\
3 & Exponential & 96 \\
\bottomrule
\end{tabular}
\end{center}

\begin{enumerate}
    \item What grade does the decaying average give this student?
    \vspace{1.5cm}
    
    \item Is the student proficient in all topics?
    \vspace{1.5cm}
    
    \item How might the system interpret this pattern?
    \vspace{1.5cm}
\end{enumerate}

\section*{Part 4: Detecting the Limitation}

Discuss:

\begin{enumerate}
    \item What assumption does the decaying average make about learning?
    \vspace{1.5cm}
    
    \item Why might this fail when topics change?
    \vspace{1.5cm}
    
    \item What situations does this system model well?
    \vspace{1.5cm}
    
    \item When does it fail?
    \vspace{1.5cm}
\end{enumerate}

\section*{Part 5: Designing a Better System}

You are now a grading system designer.

Propose guidelines for teachers:

\begin{itemize}
    \item How should assignments be structured?
    \item Should topics be revisited?
    \item How many data points are needed per skill?
\end{itemize}

\vspace{3cm}

\section*{Final Report}

Write a report explaining:

\begin{enumerate}
    \item How the decaying average works
    \item When it measures learning well
    \item When it fails
    \item Your recommended best practices
\end{enumerate}

Support your arguments with data and reasoning.

\end{document}